% Partie : histoire
% Chapitre : biographies
% Section : Rózsa Péter

\documentclass[../../../include/open-logic-section]{subfiles}

\begin{document}

\olfileid[fr]{his}{bio}{pet}

\olsection{R\'ozsa P\'eter}

\olphoto{peter-rozsa}{R\'ozsa P\'eter}

R\'ozsa P\'eter
naquit sous
le nom de
R\'ozsa Politzer
à Budapest,
en Hongrie,
le 17 février 1905.
Elle est surtout connue pour ses travaux fondateurs
sur les fonctions récursives.

P\'eter grandit
dans une période
politique difficile :
la Première Guerre
mondiale faisait
rage pendant
son adolescence.
Elle put néanmoins
fréquenter la
prestigieuse école
de jeunes filles
Maria Terezia
à Budapest,
où elle termina
ses études
en 1922.
Elle étudia
ensuite à
l'université
P\'azm\'any P\'eter
de Budapest,
devenue plus tard
l'université
E\"otv\"os Lor\'and.
À la demande
de son père,
elle commença
par la chimie,
puis se tourna
vers les
mathématiques
et obtint
son diplôme
en 1927.
Elle avait
les qualifications
pour enseigner
les mathématiques
au secondaire,
mais la situation
économique était
très difficile,
la Grande Dépression
touchant l'économie
mondiale. Elle
gagna alors
sa vie grâce
à divers emplois
de tutorat
et de cours
particuliers
de mathématiques.
Elle retourna
finalement à
l'université
pour poursuivre
des études
supérieures
en mathématiques.

P\'eter avait
d'abord prévu
de travailler
en théorie
des nombres.
En découvrant
que ses résultats
avaient déjà
été démontrés,
elle faillit
abandonner
les mathématiques.
On l'encouragea
à étudier les
théorèmes
d'incomplétude
de G\"odel,
et elle redémontra
sans le savoir
plusieurs de
ses résultats
par d'autres
méthodes.
Cette expérience
lui rendit
confiance.
Inspirée par
le programme
fondationnel
de David Hilbert,
elle rédigea
ses premiers articles sur la théorie
des fonctions récursives.
Elle obtint
son doctorat
en 1935 et
devint, en 1937,
rédactrice du
\emph{Journal of
Symbolic Logic}.

Ses premiers
articles sont
largement reconnus
comme fondateurs
pour la théorie
des fonctions
récursives.
Dans
\cite{Peter1935a},
elle étudia
les liens
entre différents
types de récursion.
Dans
\cite{Peter1935b},
elle montra
qu'une certaine
fonction définie
par récursion
n'est pas
récursive primitive.
Elle simplifiait
ainsi un résultat
antérieur de
Wilhelm Ackermann.
Cette fonction
simplifiée est
aujourd'hui
souvent appelée
fonction d'Ackermann,
et parfois,
plus justement,
fonction
d'Ackermann--P\'eter.
Elle écrivit
le premier
livre consacré
à la théorie
des fonctions
récursives
\citep{Peter1951}.

Malgré l'importance
et l'influence
de son travail,
P\'eter n'obtint
un poste d'enseignement
à temps plein
qu'en 1945.
Pendant l'occupation
nazie de la
Hongrie, lors
de la Seconde
Guerre mondiale,
les lois antisémites
l'empêchèrent
d'enseigner.
En 1944,
le gouvernement
créa un ghetto
juif à Budapest,
isolé du reste
de la ville
et surveillé
par des gardes
armés. P\'eter
fut contrainte
d'y vivre
jusqu'à sa
libération
en 1945.
Elle enseigna
ensuite à
l'École de
formation des
enseignants
de Budapest,
puis, à partir
de 1955,
à l'université
E\"otv\"os Lor\'and.
Elle fut la
première mathématicienne
hongroise à
obtenir le
titre de
docteure de
l'Académie en
mathématiques,
et fut élue,
en 1973,
membre correspondante
de l'Académie hongroise
des sciences.

P\'eter était
connue pour
sa passion
de l'enseignement
des mathématiques.
Elle préférait
explorer avec
ses étudiants
la nature
et la beauté
des problèmes
plutôt que
de se limiter
à des cours
magistraux.
Ses étudiants
l'appelaient
affectueusement
«~Tante Rosa~».
Elle mourut
en 1977
à l'âge
de~71 ans.

\begin{reading}
Pour d'autres
textes biographiques,
voir
\citep{Oconnor2014}
et
\citep{Andrasfai1986}.
\citet{Tamassy1994}
a réalisé
un bref entretien
avec P\'eter.
Pour une lecture
plaisante sur
les mathématiques,
voir son livre
\emph{Playing
With Infinity}
\citep{Peter2010}.
\end{reading}

\end{document}
